\documentclass[11pt]{article}

\usepackage[utf8]{inputenc}
\usepackage[T1]{fontenc}
\usepackage{lmodern}
\usepackage{geometry}
\usepackage{amsmath,amssymb,amsfonts,amsthm,mathtools}
\usepackage{bm}
\usepackage{enumitem}
\usepackage{microtype}
\usepackage{hyperref}
\usepackage{titlesec}
\usepackage{booktabs}
\usepackage{array}
\usepackage{setspace}
\usepackage{mathrsfs}
\usepackage{physics}

\geometry{margin=1in}
\hypersetup{colorlinks=true, linkcolor=black, urlcolor=blue, citecolor=black}
\setlist[enumerate]{topsep=0.4em,itemsep=0.35em}
\setlist[itemize]{topsep=0.3em,itemsep=0.25em}
\titleformat{\section}{\large\bfseries}{\thesection.}{0.5em}{}
\titleformat{\subsection}{\normalsize\bfseries}{\thesubsection.}{0.5em}{}
\setstretch{1.06}

\newcommand{\Aether}{\AE{}ther}
\newcommand{\Aflow}{\AE{}ther-flow}
\newcommand{\TheoryName}{\Aflow{} Interpretation of Relativity}
\newcommand{\TheoryProgram}{\Aether{} / \Aflow{} framework}

\newcommand{\CompletedMetric}{g^{\sharp}}

\newcommand{\Car}{\mathrm{car}}
\newcommand{\Cur}{\mathrm{cur}}
\newcommand{\Hom}{\mathrm{hom}}

\newcommand{\ForSpace}[1]{\mathcal I_{#1}^{\mathrm{for}}}
\newcommand{\PrimShell}{\mathcal S_{\mathrm{prim}}}

\newcommand{\Reduction}[1]{\mathcal R_{#1}}
\newcommand{\CurrentCarrier}[1]{\mathfrak C_{#1}^{\mathrm{cur}}}
\newcommand{\EqCarrier}[1]{\mathfrak C_{#1}^{\mathrm{eq}}}
\newcommand{\CurrentExtract}[1]{\mathscr E_{#1}^{\mathrm{cur}}}
\newcommand{\EqExtract}[1]{\mathscr E_{#1}^{\mathrm{eq}}}
\newcommand{\PrimCurrent}[1]{\mathcal J_{#1}^{\mathrm{prim}}}
\newcommand{\PrimTheta}[1]{\Theta_{#1}^{\mathrm{prim}}}

\newcommand{\MetricOut}{\mathscr G_{\mathrm{out}}}
\newcommand{\MatterOut}{\mathscr T_{\mathrm{out}}}
\newcommand{\DeepMetricOut}{\mathscr G_{\mathrm{deep}}}
\newcommand{\DeepMatterOut}{\mathscr T_{\mathrm{deep}}}

\newcommand{\SubSpace}[1]{\mathcal M_{#1}}
\newcommand{\DataProj}[1]{\pi_{#1}}
\newcommand{\SubVisible}[1]{\pi_{#1}^{X}}
\newcommand{\SubCurrent}[1]{\pi_{#1}^{\mathrm{cur}}}
\newcommand{\SubEq}[1]{\pi_{#1}^{\mathrm{eq}}}
\newcommand{\Fiber}[1]{\mathcal F_{#1}}
\newcommand{\ShellEmbed}[1]{\iota_{#1}^{\mathrm{sh}}}
\newcommand{\StationaryFunctional}[1]{\mathscr A_{#1}}

\newcommand{\DeepSpace}[1]{\mathcal U_{#1}}
\newcommand{\VisibleMap}[1]{\mathcal V_{#1}}
\newcommand{\DeepCurrentCarrier}[1]{\mathfrak C_{#1}^{\mathrm{deep,cur}}}
\newcommand{\DeepEqCarrier}[1]{\mathfrak C_{#1}^{\mathrm{deep,eq}}}
\newcommand{\ChainSelect}[1]{\mathcal S_{#1}}
\newcommand{\DeepCurrent}[1]{\mathcal J_{#1}^{\mathrm{deep}}}
\newcommand{\DeepTheta}[1]{\Theta_{#1}^{\mathrm{deep}}}
\newcommand{\DeepShell}{\mathcal U_{\mathrm{tot}}}
\newcommand{\ChainSelectTriple}{\mathcal S_{\mathrm{tot}}}

\newtheorem{definition}{Definition}
\newtheorem{lemma}{Lemma}
\newtheorem{proposition}{Proposition}
\newtheorem{remark}{Remark}
\newtheorem{corollary}{Corollary}
\newtheorem{theorem}{Theorem}

\title{\textbf{The \TheoryName{}}\\[0.4em]
\large Primitive-Reservoir Observer-Reduction Selector\\
\large Existence and Uniqueness from Stationary Substrate Structure}
\author{}
\date{}

\begin{document}

\maketitle

\begin{abstract}
The current primitive-reservoir observer-reduction chain forcing theorem moved the active burden one honest layer lower, but it still left one main assumption in place: the existence of a chain-faithful deeper substrate package carrying sectorwise deeper spaces, visible maps, deeper current and equilibrium carriers, and selectors into the recorded primitive forcing shells. The next benchmark-facing gain is therefore not another same-output relay. It is to derive that package itself from explicit substrate structure.

The present note records the cleanest immediate step of that kind. For each sector it introduces an explicit stationary substrate construction: a substrate state space, a projection to visible/current/equilibrium data, an embedded realization of the primitive forcing shell inside that substrate space, and a nonnegative fiberwise stationary functional whose restriction to each data fiber is coercive and strictly convex. The deeper sector space is then no longer assumed. It is the image of the substrate data projection. The visible map and deeper current/equilibrium carriers are no longer separately postulated. They are the coordinate projections on that image. The sectorwise forcing selector is no longer assumed either. It is the unique shell point whose embedded representative is the fiberwise stationary zero of the explicit substrate functional.

The positive result stays disciplined. This note does not claim a final microphysical action or a completed first-principles derivation of GR. Its narrower theorem is that the main remaining assumption in the chain-forcing note is removed at this level: the deeper package there is realized from explicit stationary substrate structure, the chain-separation property follows from fiberwise uniqueness, and the same one operative metric \(\CompletedMetric\) together with the same ordinary matter route are preserved without introducing a second metric or an enlarged low-energy matter law.
\end{abstract}

\section{Introduction}

The active derivational program of \TheoryName{} has already crossed the point at which another note that merely relays the same primitive observer-reduction chain would count as a primary gain. The primitive current and equilibrium package is already derived. The observer-sector outputs already factor through the primitive defect fibers. The one operative metric output and the ordinary matter route are already fixed on that same chain. The later forcing theorem then showed that, if a chain-faithful deeper substrate package exists, the primitive chain is itself forced by deeper explicit structure while preserving the same benchmark one-metric package.

That theorem changes the burden sharply. Once the forcing note is on record, the remaining assumption is no longer the acceptability of the primitive chain itself. It is the existence of the deeper package that the forcing theorem used. The next honest benchmark-facing move is therefore to explain where the deeper sector spaces, visible maps, deeper current/equilibrium carriers, and shell selectors come from, and why the chain-separation property needed for uniqueness holds rather than being separately stipulated.

The present note attacks exactly that narrower burden and no more. It does not claim that the final substrate microphysics has already been found. It does not enlarge the low-energy matter law, and it does not alter the benchmark exact-closure package. Its positive result is only that, once one works with an explicit sectorwise stationary substrate construction, the deeper package required by the forcing theorem is no longer an extra assumption. It is generated canonically from the stationary substrate structure itself.

\paragraph{Framework and claim boundary.}
\TheoryProgram{} denotes the broader \Aether{} / \Aflow{} framework built on the claims that the \Aether{} is the underlying four-dimensional substrate of reality, the \Aflow{} is its intrinsic ordered motion, observed three-dimensional space is the local experiential slice of that deeper substrate, S-time is the experienced order of change arising from matter, light, and the \Aflow{}, and observed expansion is the three-dimensional appearance of deeper four-dimensional ordered motion. Within that framework, \TheoryName{} denotes the exact-closure theory in which the effective gravitational dynamics are adopted to be exactly Einsteinian with universal matter coupling. In the active sequence, \emph{adoption} means use of that established relativistic dynamics without claiming substrate derivation, while \emph{derivation} is reserved for a first-principles recovery from explicit substrate variables.

The present note stays strictly inside that boundary. It does not claim a completed first-principles derivation of GR from final substrate equations. Its narrower theorem is only that the selector-existence / selector-uniqueness mechanism required below the current chain-forcing note can be obtained from explicit stationary substrate structure while preserving the same benchmark one-metric package.

\section{Fixed Primitive Continuation Data}

\begin{definition}[Recorded primitive continuation data]
\label{def:fixed}
Fix the following continuation data from the active primitive-reservoir observer-reduction line.
\begin{enumerate}
    \item For each sector label \(\sigma\in\{\Car,\Cur,\Hom\}\), a primitive forcing shell
    \[
    \ForSpace{\sigma},
    \]
    a visible reduction map
    \[
    \Reduction{\sigma}:\ForSpace{\sigma}\to X_{\sigma},
    \]
    and shell-level current and equilibrium carriers
    \[
    \CurrentCarrier{\sigma}:\ForSpace{\sigma}\to Z_{\sigma}^{\mathrm{cur}},
    \qquad
    \EqCarrier{\sigma}:\ForSpace{\sigma}\to Z_{\sigma}^{\mathrm{eq}}.
    \]
    \item The product shell
    \[
    \PrimShell
    \equiv
    \ForSpace{\Car}\times \ForSpace{\Cur}\times \ForSpace{\Hom}.
    \]
    \item Fixed shell extractors
    \[
    \CurrentExtract{\sigma}:Z_{\sigma}^{\mathrm{cur}}\to Y_{\sigma}^{\mathrm{cur}},
    \qquad
    \EqExtract{\sigma}:Z_{\sigma}^{\mathrm{eq}}\to Y_{\sigma}^{\mathrm{eq}},
    \]
    together with the recorded shell composites
    \begin{equation}
    \PrimCurrent{\sigma}
    =
    \CurrentExtract{\sigma}\circ \CurrentCarrier{\sigma},
    \qquad
    \PrimTheta{\sigma}
    =
    \EqExtract{\sigma}\circ \EqCarrier{\sigma}.
    \label{eq:primcomposites}
    \end{equation}
    \item The recorded metric and ordinary-matter outputs
    \[
    \MetricOut:\PrimShell\to\{\text{observer metrics}\},
    \qquad
    \MatterOut:\PrimShell\times\Psi\to\{\text{observer stress tensors}\},
    \]
    satisfying
    \begin{equation}
    \MetricOut(\mathbf w)=\CompletedMetric,
    \qquad
    \MatterOut(\mathbf w,\psi)
    =
    T_{\mu\nu}^{\mathrm{matter}}[\CompletedMetric,\psi]
    \label{eq:fixedoutputs}
    \end{equation}
    for every \(\mathbf w\in\PrimShell\) and every matter configuration \(\psi\in\Psi\).
\end{enumerate}
\end{definition}

\begin{remark}
The present note does not reopen the primitive shell or the benchmark output maps. Those remain fixed continuation data. The point here is narrower: derive from still more explicit substrate structure the deeper sector spaces and selectors that the current chain-forcing theorem had to assume.
\end{remark}

\section{Explicit Stationary Substrate Construction}

\begin{definition}[Sectorwise explicit stationary substrate construction]
\label{def:construction}
Fix \(\sigma\in\{\Car,\Cur,\Hom\}\). An \emph{explicit stationary substrate construction} for sector \(\sigma\) consists of the following data.
\begin{enumerate}
    \item A substrate state manifold
    \[
    \SubSpace{\sigma}
    \]
    together with a smooth data projection
    \[
    \DataProj{\sigma}:\SubSpace{\sigma}\to
    X_{\sigma}\times Z_{\sigma}^{\mathrm{cur}}\times Z_{\sigma}^{\mathrm{eq}}.
    \]
    Write its component maps as
    \[
    \SubVisible{\sigma},\qquad
    \SubCurrent{\sigma},\qquad
    \SubEq{\sigma},
    \]
    so that
    \[
    \DataProj{\sigma}(m)
    =
    \bigl(
    \SubVisible{\sigma}(m),
    \SubCurrent{\sigma}(m),
    \SubEq{\sigma}(m)
    \bigr).
    \]
    Define the sectorwise deeper data space by
    \[
    \DeepSpace{\sigma}
    \equiv
    \operatorname{im}\DataProj{\sigma}
    \subset
    X_{\sigma}\times Z_{\sigma}^{\mathrm{cur}}\times Z_{\sigma}^{\mathrm{eq}}.
    \]
    \item A smooth embedding of the primitive shell into the substrate state space,
    \[
    \ShellEmbed{\sigma}:\ForSpace{\sigma}\hookrightarrow \SubSpace{\sigma},
    \]
    satisfying the compatibility identities
    \begin{equation}
    \SubVisible{\sigma}\circ \ShellEmbed{\sigma}
    =
    \Reduction{\sigma},
    \qquad
    \SubCurrent{\sigma}\circ \ShellEmbed{\sigma}
    =
    \CurrentCarrier{\sigma},
    \qquad
    \SubEq{\sigma}\circ \ShellEmbed{\sigma}
    =
    \EqCarrier{\sigma}.
    \label{eq:embedcompat}
    \end{equation}
    \item A nonnegative substrate functional
    \[
    \StationaryFunctional{\sigma}:\SubSpace{\sigma}\to[0,\infty)
    \]
    which is \(C^{1}\) along the fibers of \(\DataProj{\sigma}\).
    \item For each deeper datum \(u\in\DeepSpace{\sigma}\), the fiber
    \[
    \Fiber{\sigma}(u)
    \equiv
    \DataProj{\sigma}^{-1}(u)
    \]
    is nonempty, closed, and geodesically convex in the current fiber geometry, and the restricted functional
    \[
    \StationaryFunctional{\sigma}\big|_{\Fiber{\sigma}(u)}
    \]
    is coercive, lower semicontinuous, and strictly geodesically convex in the sense that for every nonconstant geodesic
    \[
    \gamma:[0,1]\to \Fiber{\sigma}(u)
    \]
    one has
    \begin{equation}
    \StationaryFunctional{\sigma}(\gamma(t))
    <
    (1-t)\StationaryFunctional{\sigma}(\gamma(0))
    +
    t\,\StationaryFunctional{\sigma}(\gamma(1)),
    \qquad
    0<t<1.
    \label{eq:strictconvex}
    \end{equation}
    \item Exact shell realization: for every \(u\in\DeepSpace{\sigma}\),
    \begin{equation}
    \inf_{m\in\Fiber{\sigma}(u)}
    \StationaryFunctional{\sigma}(m)
    =0,
    \label{eq:zeroinf}
    \end{equation}
    and a point \(m\in\SubSpace{\sigma}\) satisfies
    \[
    \StationaryFunctional{\sigma}(m)=0
    \]
    if and only if \(m\in\operatorname{im}\ShellEmbed{\sigma}\).
\end{enumerate}
Define the induced deeper visible map and deeper carriers by coordinate restriction:
\begin{equation}
\VisibleMap{\sigma}(x,z^{\mathrm{cur}},z^{\mathrm{eq}})\equiv x,
\quad
\DeepCurrentCarrier{\sigma}(x,z^{\mathrm{cur}},z^{\mathrm{eq}})
\equiv z^{\mathrm{cur}},
\quad
\DeepEqCarrier{\sigma}(x,z^{\mathrm{cur}},z^{\mathrm{eq}})
\equiv z^{\mathrm{eq}}.
\label{eq:deepcoords}
\end{equation}
\end{definition}

\begin{remark}
Definition \ref{def:construction} removes the first layer of arbitrariness from the current forcing theorem. The deeper sector space is now the image of an explicit substrate projection, not a separately introduced set. The deeper visible map and deeper current/equilibrium carriers are now coordinate projections on that image, not separate postulates.
\end{remark}

\section{Existence and Uniqueness of the Sectorwise Forcing Selectors}

\begin{theorem}[Unique fiberwise stationary representatives]
\label{thm:minimizer}
Fix \(\sigma\in\{\Car,\Cur,\Hom\}\), and assume Definitions \ref{def:fixed} and \ref{def:construction}. Then for every \(u\in\DeepSpace{\sigma}\), there exists a unique point
\[
m_{\sigma}(u)\in\Fiber{\sigma}(u)
\]
such that
\[
\StationaryFunctional{\sigma}(m_{\sigma}(u))=0.
\]
This point is the unique minimizer of \(\StationaryFunctional{\sigma}\) on \(\Fiber{\sigma}(u)\), and it is stationary for the fiber-restricted variational problem.
\end{theorem}

\begin{proof}
Fix \(u\in\DeepSpace{\sigma}\). By Definition \ref{def:construction}, the fiber \(\Fiber{\sigma}(u)\) is nonempty, and the restricted functional is coercive and lower semicontinuous there. The direct method of the calculus of variations therefore yields at least one minimizer \(m_{\sigma}(u)\) of \(\StationaryFunctional{\sigma}\) on \(\Fiber{\sigma}(u)\). By \eqref{eq:zeroinf}, the minimum value is \(0\), so
\[
\StationaryFunctional{\sigma}(m_{\sigma}(u))=0.
\]

To prove uniqueness, suppose \(m,m'\in\Fiber{\sigma}(u)\) both satisfy
\[
\StationaryFunctional{\sigma}(m)=
\StationaryFunctional{\sigma}(m')=0.
\]
If \(m\neq m'\), geodesic convexity of the fiber gives a nonconstant geodesic \(\gamma\) joining them inside \(\Fiber{\sigma}(u)\). Strict convexity \eqref{eq:strictconvex} then gives, for \(0<t<1\),
\[
\StationaryFunctional{\sigma}(\gamma(t))
<
(1-t)\StationaryFunctional{\sigma}(m)
+
t\,\StationaryFunctional{\sigma}(m')
=
0,
\]
contradicting nonnegativity of \(\StationaryFunctional{\sigma}\). Hence the zero point is unique. Because the functional is \(C^{1}\) along fibers, its unique fiberwise minimizer is stationary for the fiber-restricted variational problem.
\end{proof}

\begin{definition}[Sectorwise forcing selectors]
\label{def:selector}
For each \(\sigma\in\{\Car,\Cur,\Hom\}\), define the selector
\[
\ChainSelect{\sigma}:\DeepSpace{\sigma}\to\ForSpace{\sigma}
\]
by the condition
\begin{equation}
\ShellEmbed{\sigma}\bigl(\ChainSelect{\sigma}(u)\bigr)
=
m_{\sigma}(u),
\label{eq:selectordef}
\end{equation}
where \(m_{\sigma}(u)\) is the unique fiberwise stationary zero from Theorem \ref{thm:minimizer}.
\end{definition}

\begin{proposition}[Selectors satisfy the chain identities]
\label{prop:compat}
For every \(\sigma\in\{\Car,\Cur,\Hom\}\),
\begin{equation}
\Reduction{\sigma}\circ\ChainSelect{\sigma}
=
\VisibleMap{\sigma},
\label{eq:visiblecompat}
\end{equation}
\begin{equation}
\CurrentCarrier{\sigma}\circ\ChainSelect{\sigma}
=
\DeepCurrentCarrier{\sigma},
\qquad
\EqCarrier{\sigma}\circ\ChainSelect{\sigma}
=
\DeepEqCarrier{\sigma}.
\label{eq:carriercompat}
\end{equation}
\end{proposition}

\begin{proof}
Fix \(u\in\DeepSpace{\sigma}\). By \eqref{eq:selectordef},
\[
\ShellEmbed{\sigma}\bigl(\ChainSelect{\sigma}(u)\bigr)=m_{\sigma}(u).
\]
Applying the three identities in \eqref{eq:embedcompat} gives
\[
\Reduction{\sigma}\bigl(\ChainSelect{\sigma}(u)\bigr)
=
\SubVisible{\sigma}\bigl(m_{\sigma}(u)\bigr),
\]
\[
\CurrentCarrier{\sigma}\bigl(\ChainSelect{\sigma}(u)\bigr)
=
\SubCurrent{\sigma}\bigl(m_{\sigma}(u)\bigr),
\qquad
\EqCarrier{\sigma}\bigl(\ChainSelect{\sigma}(u)\bigr)
=
\SubEq{\sigma}\bigl(m_{\sigma}(u)\bigr).
\]
But \(\DataProj{\sigma}(m_{\sigma}(u))=u\), and the maps in \eqref{eq:deepcoords} are exactly the coordinate projections on \(\DeepSpace{\sigma}\). Hence
\[
\SubVisible{\sigma}\bigl(m_{\sigma}(u)\bigr)=\VisibleMap{\sigma}(u),
\]
\[
\SubCurrent{\sigma}\bigl(m_{\sigma}(u)\bigr)=\DeepCurrentCarrier{\sigma}(u),
\qquad
\SubEq{\sigma}\bigl(m_{\sigma}(u)\bigr)=\DeepEqCarrier{\sigma}(u),
\]
which proves \eqref{eq:visiblecompat} and \eqref{eq:carriercompat}.
\end{proof}

\begin{theorem}[Chain-separation is forced by the stationary construction]
\label{thm:separation}
Fix \(\sigma\in\{\Car,\Cur,\Hom\}\). If \(w,w'\in\ForSpace{\sigma}\) satisfy
\begin{equation}
\Reduction{\sigma}(w)=\Reduction{\sigma}(w'),
\qquad
\CurrentCarrier{\sigma}(w)=\CurrentCarrier{\sigma}(w'),
\qquad
\EqCarrier{\sigma}(w)=\EqCarrier{\sigma}(w'),
\label{eq:samedata}
\end{equation}
then \(w=w'\).
\end{theorem}

\begin{proof}
Assume \eqref{eq:samedata}. By the compatibility identities \eqref{eq:embedcompat},
\[
\DataProj{\sigma}\bigl(\ShellEmbed{\sigma}(w)\bigr)
=
\DataProj{\sigma}\bigl(\ShellEmbed{\sigma}(w')\bigr).
\]
Call the common value \(u\in\DeepSpace{\sigma}\). Because \(\ShellEmbed{\sigma}(w)\) and \(\ShellEmbed{\sigma}(w')\) lie in the image of the primitive shell embedding, Definition \ref{def:construction} gives
\[
\StationaryFunctional{\sigma}\bigl(\ShellEmbed{\sigma}(w)\bigr)=0,
\qquad
\StationaryFunctional{\sigma}\bigl(\ShellEmbed{\sigma}(w')\bigr)=0.
\]
Thus both points are stationary zeros in the same fiber \(\Fiber{\sigma}(u)\). By Theorem \ref{thm:minimizer}, that fiber has a unique zero point. Therefore
\[
\ShellEmbed{\sigma}(w)=\ShellEmbed{\sigma}(w').
\]
Since \(\ShellEmbed{\sigma}\) is an embedding, it is injective, so \(w=w'\).
\end{proof}

\begin{corollary}[Total selector exists uniquely]
\label{cor:totalselector}
Define
\[
\DeepShell
\equiv
\DeepSpace{\Car}\times\DeepSpace{\Cur}\times\DeepSpace{\Hom},
\]
and
\[
\ChainSelectTriple(u_{\Car},u_{\Cur},u_{\Hom})
\equiv
\bigl(
\ChainSelect{\Car}(u_{\Car}),
\ChainSelect{\Cur}(u_{\Cur}),
\ChainSelect{\Hom}(u_{\Hom})
\bigr).
\]
Then \(\ChainSelectTriple\) is well defined and uniquely determined by the explicit stationary substrate construction.
\end{corollary}

\begin{proof}
Theorem \ref{thm:minimizer} and Definition \ref{def:selector} give a unique selector in each sector. Their product is therefore a well-defined unique total selector.
\end{proof}

\section{Induced Primitive Readouts and the Same One-Metric Output}

\begin{proposition}[The primitive current and equilibrium composites are induced]
\label{prop:readouts}
For each \(\sigma\in\{\Car,\Cur,\Hom\}\), define
\[
\DeepCurrent{\sigma}
\equiv
\CurrentExtract{\sigma}\circ \DeepCurrentCarrier{\sigma},
\qquad
\DeepTheta{\sigma}
\equiv
\EqExtract{\sigma}\circ \DeepEqCarrier{\sigma}.
\]
Then
\begin{equation}
\DeepCurrent{\sigma}
=
\PrimCurrent{\sigma}\circ \ChainSelect{\sigma},
\qquad
\DeepTheta{\sigma}
=
\PrimTheta{\sigma}\circ \ChainSelect{\sigma}.
\label{eq:readoutpullback}
\end{equation}
\end{proposition}

\begin{proof}
Using \eqref{eq:primcomposites} and the carrier compatibilities \eqref{eq:carriercompat},
\[
\PrimCurrent{\sigma}\circ \ChainSelect{\sigma}
=
\CurrentExtract{\sigma}\circ \CurrentCarrier{\sigma}\circ \ChainSelect{\sigma}
=
\CurrentExtract{\sigma}\circ \DeepCurrentCarrier{\sigma}
=
\DeepCurrent{\sigma}.
\]
The proof for \(\DeepTheta{\sigma}\) is identical.
\end{proof}

\begin{definition}[Deeper metric and ordinary-matter outputs]
\label{def:deepoutputs}
Define
\[
\DeepMetricOut:\DeepShell\to\{\text{observer metrics}\}
\]
and
\[
\DeepMatterOut:\DeepShell\times\Psi\to\{\text{observer stress tensors}\}
\]
by
\begin{equation}
\DeepMetricOut
\equiv
\MetricOut\circ \ChainSelectTriple
\label{eq:deepmetricdef}
\end{equation}
and
\begin{equation}
\DeepMatterOut(\mathbf u,\psi)
\equiv
\MatterOut(\ChainSelectTriple(\mathbf u),\psi).
\label{eq:deepmatterdef}
\end{equation}
\end{definition}

\begin{theorem}[The stationary construction preserves the same one operative metric and matter route]
\label{thm:deepmetric}
For every \(\mathbf u\in\DeepShell\) and every \(\psi\in\Psi\),
\begin{equation}
\DeepMetricOut(\mathbf u)=\CompletedMetric
\label{eq:deepmetric}
\end{equation}
and
\begin{equation}
\DeepMatterOut(\mathbf u,\psi)
=
T_{\mu\nu}^{\mathrm{matter}}[\CompletedMetric,\psi].
\label{eq:deepmatter}
\end{equation}
Consequently the explicit stationary substrate construction introduces neither a second operative metric nor an enlarged low-energy matter law.
\end{theorem}

\begin{proof}
Equation \eqref{eq:deepmetric} follows immediately from \eqref{eq:deepmetricdef} and \eqref{eq:fixedoutputs}:
\[
\DeepMetricOut(\mathbf u)
=
\MetricOut(\ChainSelectTriple(\mathbf u))
=
\CompletedMetric.
\]
Likewise,
\[
\DeepMatterOut(\mathbf u,\psi)
=
\MatterOut(\ChainSelectTriple(\mathbf u),\psi)
=
T_{\mu\nu}^{\mathrm{matter}}[\CompletedMetric,\psi],
\]
which proves \eqref{eq:deepmatter}. Since the deeper outputs are defined as pullbacks of the already fixed primitive outputs, they introduce no additional operative metric and no enlarged matter law.
\end{proof}

\section{Main Theorem}

\begin{theorem}[Selector existence and uniqueness from stationary substrate structure]
\label{thm:main}
Assume the fixed primitive continuation data of Definition \ref{def:fixed} and, for each sector \(\sigma\in\{\Car,\Cur,\Hom\}\), an explicit stationary substrate construction in the sense of Definition \ref{def:construction}. Then:
\begin{enumerate}
    \item the deeper sector spaces \(\DeepSpace{\sigma}\) arise canonically as the images of the substrate data projections \(\DataProj{\sigma}\);
    \item the deeper visible maps \(\VisibleMap{\sigma}\) and deeper current/equilibrium carriers \(\DeepCurrentCarrier{\sigma},\DeepEqCarrier{\sigma}\) arise canonically as the coordinate projections on those deeper sector spaces;
    \item the sectorwise forcing selectors \(\ChainSelect{\sigma}\) and the total selector \(\ChainSelectTriple\) exist and are unique;
    \item the primitive chain-separation property is forced by the stationary construction rather than assumed independently;
    \item the primitive current and equilibrium composites are induced from the deeper current and equilibrium carriers through the already fixed shell extractors;
    \item the deeper metric output is exactly the same one operative metric \(\CompletedMetric\);
    \item the deeper ordinary-matter output is exactly the same standard matter route through \(\CompletedMetric\);
    \item consequently the stationary construction introduces neither a second operative metric nor an enlarged low-energy matter law.
\end{enumerate}
Therefore the main remaining assumption of the current primitive-reservoir observer-reduction chain forcing theorem is removed at this level: the chain-faithful deeper package there is no longer separately postulated, but is realized by explicit stationary substrate structure.
\end{theorem}

\begin{proof}
Items (1) and (2) are part of Definition \ref{def:construction}. Item (3) is Theorem \ref{thm:minimizer}, Definition \ref{def:selector}, and Corollary \ref{cor:totalselector}. Item (4) is Theorem \ref{thm:separation}. Item (5) is Proposition \ref{prop:readouts}. Items (6)--(8) are Theorem \ref{thm:deepmetric}. The final sentence is exactly the routing consequence of those items taken together.
\end{proof}

\begin{corollary}[What remains open after the selector theorem]
\label{cor:next}
After Theorem \ref{thm:main}, the remaining honest burden is no longer the existence or uniqueness of the sectorwise selectors once explicit stationary substrate structure is granted. The remaining burden is to derive or justify that stationary substrate structure itself --- including the projected substrate spaces, the coercive and strictly convex fiberwise functionals, and the exact shell realization property --- from yet more explicit substrate equations, symmetries, or constraints while preserving the same benchmark one-metric package.
\end{corollary}

\begin{proof}
Theorem \ref{thm:main} already converts the selector mechanism from a hypothesis into an output of the stationary substrate construction. What it does not yet do is derive why the particular stationary substrate construction of Definition \ref{def:construction} should hold from a still more primitive substrate theory. That is therefore the next remaining burden.
\end{proof}

\section{Conclusion}

The positive gain of this note is focused but real. The deeper sector spaces, deeper visible maps, deeper current/equilibrium carriers, and sectorwise forcing selectors required below the current primitive observer-reduction chain are no longer simply assumed once explicit stationary substrate structure is supplied. They arise canonically from a substrate state space with a projection to visible/current/equilibrium data and a nonnegative fiberwise stationary functional whose exact zero set is the embedded primitive shell and whose coercive strict convexity forces unique fiberwise stationary zeros. In particular, the selector existence and selector uniqueness statements are now derivational outputs rather than separate hypotheses, and the chain-separation property follows from the same construction.

The benchmark boundary remains fixed. Pulling the already recorded primitive outputs back along the unique total selector yields exactly the same one operative metric \(\CompletedMetric\) and the same standard ordinary-matter route through that metric, with no second operative metric and no enlarged low-energy matter law. This is still not a completed first-principles derivation of GR from final substrate microphysics, and stronger promotion language should remain paused. The remaining honest burden is one layer deeper again: derive or justify the stationary substrate construction itself from yet more explicit substrate equations, symmetries, or constraint structure.

\end{document}
